\section{Evaluation}
\label{sec:evaluation}

Figure~\ref{fig:evaluation} gives an overview of the experiments.

% Figure 2: the evaluation design, on a fixed grid. Counts come from generated/numbers.tex.
\begin{figure}[t]
\centering
\hyphenpenalty=10000
\begin{tikzpicture}[
  font=\small,
  box/.style={draw, rounded corners=2pt, align=center, inner sep=3pt, minimum height=10mm,
    text width=#1, anchor=north},
  box/.default=32mm,
  data/.style={box=#1, fill=black!6},
  data/.default=32mm,
  model/.style={box=#1, fill=blue!10, draw=blue!60!black},
  model/.default=60mm,
  out/.style={box=#1, fill=orange!20},
  out/.default=60mm,
  arr/.style={-{Stealth[length=2mm]}, thick},
]
% Inputs
\node[data=44mm] (qasper) at (2.3, 0) {QASPER v0.3, development\\[1pt]\footnotesize
  \BaselinePapers{} papers, \BaselineQuestions{} questions};
\node[data] (pdfs) at (6.6, 0) {arXiv PDFs\\[1pt]\footnotesize \BaselinePdfs{} of \BaselinePapers{} downloaded};
\node[box=60mm] (app) at (11.8, 0) {The app's own code\\[1pt]\footnotesize
  \texttt{split\_pdf}, \texttt{LocalVectorStore}, \texttt{build\_messages}};
\draw[arr] (qasper.east |- pdfs.west) -- (pdfs.west);
\draw[arr] (pdfs.east) -- (pdfs.east -| app.west);

% Retrieval only (no language model)
\node[box] (sweep) at (1.85, -2.6) {Retrieval sweep\\[1pt]\footnotesize \SweepSetups{} set-ups:
  3 chunk sizes $\times$ 12 retrievers};
\node[box] (multi) at (5.35, -2.6) {Other papers loaded\\[1pt]\footnotesize $N = 0, 2, 4, 9$;
  3 retrieval rules};
\node[out=67mm] (match) at (3.6, -4.6) {Evidence in the context?\\[1pt]\footnotesize chunks
  matched to gold paragraphs by word 5-grams; no language model};
\draw[arr] (sweep.south) -- (sweep.south |- match.north);
\draw[arr] (multi.south) -- (multi.south |- match.north);

% Answers
\node[box=60mm, align=left] (conds) at (11.8, -2.6) {Answer conditions\\[1pt]\footnotesize
  shipped: system prompt only\\
  fixed: the app's top 4 chunks\\
  oracle: the annotators' evidence\\
  candidates: 500-character chunks, top 8, with and without reranking};
\node[model, below=5mm of conds] (llm) {Answer: Llama~3.1~8B in Ollama\\[1pt]\footnotesize
  temperature 0.1, seed 0};
\node[model, below=5mm of llm] (judge) {Judge: Qwen3-32B\\[1pt]\footnotesize correct, partial or
  incorrect; empty answers by rule};
\node[out, below=5mm of judge] (f1) {Judged correctness and official Answer F1};
\draw[arr] (conds) -- (llm);
\draw[arr] (llm) -- (judge);
\draw[arr] (judge) -- (f1);

% From the app's code to both branches
\coordinate (bus) at (0, -1.75);
\draw[thick] (app.south) -- (app.south |- bus);
\draw[thick] (sweep.north |- bus) -- (app.south |- bus);
\draw[arr] (sweep.north |- bus) -- (sweep.north);
\draw[arr] (multi.north |- bus) -- (multi.north);
\draw[arr] (app.south |- bus) -- (conds.north);

% Statistics
\path (f1.south) ++(0, -6mm) coordinate (statsy);
\coordinate (mid) at (7.5, 0);
\node[box=150mm] (stats) at (mid |- statsy) {Wilson 95\% intervals; exact McNemar tests paired by
  question (gained / lost), Holm-corrected within each family};
\draw[arr] (match.south) -- (match.south |- stats.north);
\draw[arr] (f1.south) -- (f1.south |- stats.north);
\end{tikzpicture}
\caption{Evaluation design. Every experiment runs the app's own retrieval and prompting code on
the arXiv PDFs of the QASPER development papers. Retrieval experiments score whether a gold
evidence paragraph reached the context, without calling a language model; answer experiments
generate with Llama~3.1~8B and grade with a Qwen3-32B judge and the official Answer F1.
Comparisons are paired by question.}
\label{fig:evaluation}
\end{figure}


\subsection{Data}

We used the development split of QASPER v0.3~\cite{dasigi2021qasper}, which contains 281 NLP
papers and 1,005 questions. Each question was answered by one or more annotators, who also
marked the paragraphs that hold the evidence. We kept the test split for a final comparison once
the design is settled. Because the app reads PDFs and not the dataset's parsed text, we
downloaded each paper from arXiv, with one request every three seconds as arXiv asks of automated
clients. 280 of the 281 downloads succeeded. The current arXiv version of a paper may differ from
the one QASPER was built from, and we measure the effect of this through evidence findability
(below).

\subsection{Pipeline under test}

All retrieval and prompting used the app's own functions (\texttt{Utils.split\_pdf},
\texttt{LocalVectorStore} and \texttt{Utils.build\_messages}). The only difference from the app
is that the model ran in Ollama instead of on Groq, behind a client with the same interface, so
that every run could be repeated on our own hardware. We used \texttt{llama3.1:8b}, which has the
same weights as Groq's \texttt{llama-3.1-8b-instant}, in Ollama's 4-bit Q4\_K\_M quantization.
This may answer slightly differently from Groq's deployment. Generation used the app's settings
(temperature 0.1 and at most 1,000 tokens), with seed 0 and a 16,384-token context window, which
is larger than any prompt in the evaluation. Each question was asked on its own, without
conversation history.

We compared three conditions. \emph{Shipped} sends what the app sent for the first question of a
conversation before the fix, which is the system prompt alone. \emph{Fixed} sends the fixed app's
prompt with the four retrieved chunks. \emph{Oracle} sends the same prompt with the annotators'
evidence paragraphs in place of the retrieved chunks. The oracle condition shows how well this
model and prompt could do if retrieval were perfect.

\subsection{Evidence retrieval}
\label{sec:ev-evidence}

We scored retrieval by whether a chunk holding a gold evidence paragraph was among the chunks
sent to the model. QASPER's paragraphs come from a parse of the paper's source, while the chunks
come from the PDF text, so the two had to be matched approximately. Both were normalised
(Unicode NFKC, line-break hyphens removed, lower case and alphanumeric tokens only). A chunk was
taken to hold a paragraph if it contained at least half of the paragraph's word 5-grams, or, for
paragraphs longer than a chunk, if at least half of the chunk's own 5-grams were in the
paragraph. We call a question's evidence \emph{findable} if some chunk of the whole PDF holds it,
and we compute retrieval rates over the questions with findable evidence. As in QASPER's
text-only evidence, figures and tables are left out.

\subsection{Retrieval sweep}

To choose retrieval settings without calling a language model, we ranked every paper's chunks for
every question under \SweepSetups{} set-ups. These covered three chunk sizes (500, 1,000 and
1,500 characters, with 20\% overlap) and several retrievers: dense retrieval with
all-MiniLM-L6-v2~\cite{wang2020minilm}, bge-small-en-v1.5~\cite{xiao2023cpack} or
e5-small-v2~\cite{wang2022e5}; BM25 keyword search~\cite{robertson2009bm25}; and hybrids of the
dense and BM25 rankings by reciprocal rank fusion~\cite{cormack2009rrf}. Five of these (MiniLM,
BM25 and the three hybrids) were also run with their top 20 chunks reranked by a cross-encoder
trained on MS~MARCO~\cite{nogueira2019rerank,bajaj2016msmarco}
(\texttt{cross-encoder/ms-marco-MiniLM-L-6-v2}). Four chunks of 1,500 characters hold three
times as much text as four chunks of 500, so we compared the set-ups at the same budget of about
4,000 characters, which is the top 8, 4 or 3 chunks depending on the chunk size. The comparison
uses the \SweepQuestions{} questions whose evidence is findable at every chunk size.

\subsection{Several papers}

To test whether other papers crowd out the one a question is about, we loaded each paper into the
app's vector store together with 2, 4 or 9 other development papers drawn at random (seed 0). We
then searched with three rules: the top four chunks; the app's rule for several files, which takes
the top four and then one chunk from each file not yet represented, up to eight; and the top four
chunks of the question's own paper, which is what happens when the user selects that paper under
``Ask about''.

\subsection{Judging answers}

We graded the answers in two ways. The official QASPER evaluator computes token F1 against each
annotator's answer and keeps the best score. It suits short extractive answers but penalises the
app's full-sentence answers, so we report it without relying on it. Our main measure was a
judgement by Qwen3-32B~\cite{yang2025qwen3} (Q4\_K\_M, reasoning enabled, temperature 0). The
judge was given the question, every annotator's answer and the app's answer, and asked whether
the app's answer gives the substance of at least one reference answer (correct), part of it
(partially correct) or neither (incorrect). The prompt is in the SI. Language-model judges have
known biases~\cite{zheng2023judge}, and in a pilot on nine answers our judge gave credit to empty
answers. We therefore graded empty answers as incorrect by rule. Three of the fixed app's
judgements came back empty because the judge's reasoning filled its 8,192-token window. These
were graded again with a 20,480-token window, and both attempts are kept.

To check the judge, we drew \JudgeCheckN{} answers of the fixed app at random (seed 0). Claude
(Anthropic) first graded them with the judge's rules, before any of the judge's verdicts on them
had been read (Section~\ref{sec:res-judge}). We then labelled the same answers ourselves, without
seeing the judge's verdicts or Claude's labels, on a page that showed the question, the reference
answers and the app's answer together with the judge's instructions (SI, Section~S6). Like the
judge, neither Claude nor we read the papers, and all grading was done against the references.

\subsection{Repeat runs}
\label{sec:eval-repeats}

Because the app answers at temperature 0.1, the same question can get a different answer in
another run. To measure how much this changes the results, we answered every question twice more
for each of the four set-ups compared in Section~\ref{sec:results} (the fixed app, the
annotators' evidence and the two leading retrieval set-ups). These runs used answer-model seeds 1
and 2 and were judged as before, and the original runs count as seed 0. Retrieval is
deterministic, so it was not repeated and the saved chunks were reused. Ollama with parallel
requests is not exactly repeatable even at a fixed seed (a rerun at seed 0 did not reproduce the
original answers), so the seeds should be read as independent runs rather than a controlled
source of randomness. To separate the judge's own variation, we also judged the seed-0 answers of
the fixed app a second time with judge seed 1. For each set-up we report the correct rate of each
run, the share of questions with the same verdict in all three runs, and the paired difference
between runs of the same set-up. This last figure shows how large a difference between set-ups
has to be before it can be told apart from run-to-run variation.

\subsection{Citation check}
\label{sec:eval-citations}

The app asks the model to cite the pages it used, and in the context each chunk is labelled with
its page (``[p.~4]''). We split each answer of the fixed app into the sentences that carry a
citation, where a sentence ending in a colon takes the list that follows it. For every citation,
we recorded whether the cited page was among the pages of the four chunks the model was given.
We then drew one cited statement at random (seed 0) from each answer whose cited pages were all
in the context, and asked the judge model whether the passages from the cited pages support it
(supported, partially supported or not supported; prompt in the SI). We labelled fifty of these
statements ourselves, without seeing the judge's verdicts, on a page that showed the question,
the statement and the cited passages exactly as the model saw them.

\subsection{Changes beyond retrieval settings}
\label{sec:eval-changes}

We tested three further changes against the app's set-up, on the same questions and with the same
model, prompt and judge.

\emph{Section-aware chunking} splits the PDF text at section headings so that no chunk spans two
sections, and starts each chunk with its section title. The headings are found from the lines of
text themselves, including numbered and Roman-numeral headings and the usual unnumbered names such
as ``Related Work''. On the \PapersDownloaded{} development PDFs, it found \SectionsCoverage{} of
QASPER's section names (\SectionsCoverageK{} of \SectionsCoverageN{}, leaving out run-in
paragraph headings). We compared it with the app's chunking at the app's settings (1,000
characters, top four) and at 500 characters with the top eight, first on retrieval and then on
answers.

\emph{Retry on a miss} first shows the model the four retrieved chunks and asks whether they hold
the answer. If the model says no, it writes a search query, the paper is searched again, and the
top four chunks of the new search replace the first ones. The prompt length therefore stays the
same.

\emph{Balanced context for comparisons} needed comparison questions, which QASPER does not have.
We built them from questions that were asked of several papers, such as ``What datasets are
used?''. For any two such papers, this gives a comparison question with an annotated answer for
each paper. Questions on baselines, datasets, evaluation metrics and English-only results gave
\ComparePairs{} pairs (30 per topic, drawn with seed 0). We compared the app's rule for several
files (the top four, then the best chunk of a missing file) with taking the top two chunks of each
paper, which gives a context of the same length. Each comparison answer was judged once for each
paper against that paper's references, and it counts as correct only if both judgements are
correct.

\subsection{Statistics}

Rates are given with Wilson 95\% intervals~\cite{wilson1927}. Conditions answered on the same
questions are compared with exact two-sided McNemar tests~\cite{mcnemar1947} on the questions
whose outcome differs, which we report as questions gained and lost. Within the retrieval sweep,
the \SweepComparisons{} comparisons with the app's set-up are Holm-corrected~\cite{holm1979}.
The sweep and the several-papers experiment are exploratory, as we designed them after the first
runs to explain what we had seen.
